\section*{第 \ref{ch:b2:speciation} 章——物种形成与大进化}

\begin{solution}{exo:b2:speciation:1}
物种是一组能够相互交配的自然种群，
与其他此类类群存在\omterm{def:b2:speciation:species}{生殖隔离}。该定义在下列情形中失效：
无性谱系（细菌、蛭形轮虫、蒲公英：
采用形态学或系统发生概念，即相似形态的聚类或
可诊断的谱系）；化石（形态学概念，因为
无法进行杂交）；从不相遇的异域种群（
生态或系统发生概念，因为相互交配
能否发生无从检验）——该定义还因易杂交的栎属植物与
环状物种而变得模糊。
\end{solution}

\begin{solution}{exo:b2:speciation:2}
在不同月份繁殖的青蛙：合子前（时间）隔离。骡：合子后
（杂种不育）。无法结合的精子：合子前（配子）隔离。
长势旺盛但\omterm{def:b2:angiosperm-reproduction:seed}{种子}不育的杂种：合子后（杂种不育，
或若缺陷出现于下一代，则为杂种衰败）。
鸣声不同的蟋蟀：合子前（行为）隔离。
\end{solution}

\begin{solution}{exo:b2:speciation:3}
异域分布：屏障将分布区一分为二（巴拿马地峡
两侧的鼓虾）。边缘分布：少数奠基者迁至分布区之外（
来自大陆祖先的达尔文雀）。同域分布：无屏障（
以山楂与苹果为宿主的苹果绕实蝇；
\omterm{def:b2:genome-diversification:polyploidy}{异源多倍体}小麦）。邻域分布：沿梯度分布，中部有\omterm{prop:b2:speciation:reinforcement}{杂交
带}（生长于矿渣堆上与其外围的禾草，
开花时间不同）。
\end{solution}

\begin{solution}{exo:b2:speciation:4}
物种 1 总是获胜；物种 2 总是获胜；稳定共存；
二者中哪一方获胜取决于初始数量。共存稳定的条件是
当每个物种都能在其环境容纳量下入侵对方时，
即每个物种对自身的抑制强于其对
另一物种的抑制（$K_1 > \alpha K_2$ 且 $K_2 > \beta K_1$）。
\end{solution}

\begin{solution}{exo:b2:speciation:5}
$\alpha K_2 = 400 < K_1 = 1000$：物种 1 可入侵；$\beta K_1 = 600
< K_2 = 800$：物种 2 可入侵。$N_1^{*} = (1000 - 400)/(1 - 0.3) =
857$，$N_2^{*} = (800 - 600)/0.7 = 286$；平衡时
物种 1 比其 $K$ 低 $\alpha N_2^{*} = 143$，物种 2 则低 $\beta
N_1^{*} = 514$。
\end{solution}

\begin{solution}{exo:b2:speciation:6}
$\alpha K_2 = 1200 > K_1$：在环境容纳量下，物种 1 无法入侵物种 2；
$\beta K_1 = 300 < K_2$：物种 2 可入侵物种 1。
物种 2 必然获胜，因为其个体对物种 1 的危害
大于物种 1 自身个体造成的危害，而其自身几乎
不受影响。取 $\alpha = \beta = 1.5$ 时，$\alpha K_2 = 1200 > 1000$
且 $\beta K_1 = 1500 > 800$：二者均无法入侵对方，
各自建立后都排斥另一方，先到的奠基者获胜。
\end{solution}

\begin{solution}{exo:b2:speciation:7}
四倍体 $4n$ 产生 $2n$ 配子；与\omterm{def:b2:meiosis-heredity:meiosis}{二倍体}的 $n$ 配子杂交后
产生三倍体 $3n$，其\omterm{def:b2:meiosis-heredity:meiosis}{减数分裂}无法使三组染色体配对，
并产生非整倍体、多数无生活力的配子。四倍体
只能与自身交配（或自交），以一步实现
隔离。它通常失败，因为它孤身一株——其潜在配偶
只有\omterm{def:b2:meiosis-heredity:meiosis}{二倍体}，而每次此类杂交都是徒劳的——
并且会被淹没；若能自交、
营养繁殖或反复产生，则有可能存续。
\end{solution}

\begin{solution}{exo:b2:speciation:8}
$0.01\times 10\times 10 = 1$ 个不相容性；若每方各为 30 个，$0.01\times
900 = 9$。隔离程度随\omterm{def:b2:cell-differentiation:differentiation}{分化}时间的平方而增长：
\omterm{def:b2:cell-differentiation:differentiation}{分化}时间长三倍的两个种群，隔离程度为
九倍——即雪球效应。
\end{solution}

\begin{solution}{exo:b2:speciation:9}
$w = 2\sqrt{8/0.05} = \qty{25}{km}$；$w = 0.3\sqrt{8/0.4} =
\qty{1.3}{km}$。\omterm{prop:b2:speciation:reinforcement}{强化}在第二种情形中更可能发生：杂种
\omterm{def:b2:population-genetics:fitness}{适合度}很低，因此拒绝另一类型配偶的雌性获益
很大，且接触带足够窄，使两种类型频繁相遇，
从而使这种拒绝行为被选择保留。
\end{solution}

\begin{solution}{exo:b2:speciation:10}
$\alpha K_2 = 450 < 500$，因此二者各自都能入侵对方，它们得以共存，
平衡时各为 $N^{*} = (500 - 450)/(1 - 0.81) = 263$，勉强共存，接近
排斥的边缘，$K$ 的微小变化即可打破平衡。若取
$\alpha = \beta = 0.4$，则 $N^{*} = (500 - 200)/(1 - 0.16) = 357$，二者均为该值。
\omterm{prop:b2:speciation:displacement}{性状替换}使每条等斜线远离另一物种的（$K/\alpha$
向外移动），因此交点远离坐标轴，
每个物种都更接近自身的 $K$：共存更为稳固，
且两者的数量都更多。
\end{solution}

\begin{solution}{exo:b2:speciation:11}
$R = h^{2}S = 0.7\times(-1) = \qty{-0.7}{mm}$：与观察结果一致。在
1977 年干旱期间，因无大型地雀存在，同样的中型地雀
被选择为具有\emph{更大}的喙，因为此时留下的食物是大\omterm{def:b2:angiosperm-reproduction:seed}{种子}；
2004 年，大\omterm{def:b2:angiosperm-reproduction:seed}{种子}被入侵者取走，
与之竞争的中型地雀则失利：
选择的方向由竞争者决定。
\end{solution}

\begin{solution}{exo:b2:speciation:12}
偶然形成：双位点模型表明，隔离是作为
因其他原因而固定的替换的副产物，或由\omterm{thm:b2:population-genetics:drift}{遗传漂变}造成，
而地理因素——屏障、奠基——提供了分隔，
使这一偶然性得以积累。由选择保留：在
两种类型相遇之处，针对\omterm{def:b2:population-genetics:fitness}{适合度}低的杂种的选择建立了\omterm{prop:b2:speciation:isolation}{合子前
隔离}（\omterm{prop:b2:speciation:reinforcement}{强化}），竞争使性状发生替换，从而
使二者能够共存。这一说法遗漏了通过
分裂选择实现的\omterm{prop:b2:speciation:modes}{同域物种形成}，在这里选择既创造又维持
物种，此外还遗漏了\omterm{def:b2:genome-diversification:polyploidy}{多倍性}，在那里偶然事件就是全部原因。
\end{solution}

\begin{solution}{pb:b2:speciation:1}
\textbf{1.} 物种 1：$N_1 + 0.9N_2 = 1200$，截距 $N_1 = 1200$
和 $N_2 = 1333$。物种 2：$N_2 + 0.3N_1 = 400$，截距 $N_2 =
400$ 和 $N_1 = 1333$。
\textbf{2.} $\beta K_1 = 360 < K_2 = 400$：物种 2 可入侵；
$\alpha K_2 = 360 < K_1 = 1200$：物种 1 可入侵。稳定
共存。
\textbf{3.} $N_1^{*} = (1200 - 360)/(1 - 0.27) = 1151$；$N_2^{*} =
(400 - 360)/0.73 = 55$。
\textbf{4.} $\mathrm{d}N_1/\mathrm{d}t = 0.5\times 1200\,(1 - 1218/1200)
= -9$ 只/年：处于环境容纳量的原住种群，
被新来者推入衰退。
\textbf{5.} $K_1 = 600$，$K_2 = 200$：$N_1^{*} = (600 - 180)/0.73 =
575$，$N_2^{*} = (200 - 180)/0.73 = 27$：两者都减半，而入侵者
仅剩 27 只，接近灭绝——干旱期正是
竞争最为激烈之时。
\textbf{6.} 入侵者是体型更大、喙更大的鸟：它
取食原住种群也利用的大\omterm{def:b2:angiosperm-reproduction:seed}{种子}，并且取食得更快，
而原住种群的小\omterm{def:b2:angiosperm-reproduction:seed}{种子}对它用处不大。
\textbf{7.} $S = 9.0 - 9.5 = \qty{-0.5}{mm}$；$R = 0.7\times(-0.5) =
\qty{-0.35}{mm}$。
\textbf{8.} $9.5 - 0.35 = \qty{9.15}{mm}$。
\textbf{9.} $N_1^{*} = (1200 - 200)/(1 - 0.15) = 1176$；$N_2^{*} =
(400 - 360)/0.85 = 47$。
\textbf{10.} 物种 1 等斜线的 $N_2$ 截距由
$K_1/0.9 = 1333$ 升至 $K_1/0.5 = 2400$：该直线远离
入侵者的等斜线转动，原住种群如今受每个入侵者的影响更小，
平衡点向原住种群的 $K$ 移动。
\textbf{11.} $2/0.35 \approx 6$ 年。之所以不会如此，是因为这种
选择只发生在干旱年份，在湿润年份会发生逆转，
（小\omterm{def:b2:angiosperm-reproduction:seed}{种子}大量存在之时，或入侵者稀少之时），而且
加性方差也会被耗尽。
\textbf{12.} 两百万年约为 $10^{6}$ 代；若
每代持续出现 \qty{0.35}{mm} 的变化，即使只持续其千分之一，
则喙长将改变数百毫米。选择很强，但并不恒定；
持续的定向变化是罕见的
例外，地质时间尺度上的净变化只是
大量相反阶段之间的微小残差。
\textbf{13.} 鸣声从父亲习得，雌性在
择偶时据此选择，喙形本身也作为线索。杂种
喙形中间，无论大\omterm{def:b2:angiosperm-reproduction:seed}{种子}还是小\omterm{def:b2:angiosperm-reproduction:seed}{种子}都
处理得不好，在干旱中最先饿死。
\textbf{14.} $w = 0.5\sqrt{8/0.2} = \qty{3.2}{km}$。
\textbf{15.} 该岛比\omterm{prop:b2:speciation:reinforcement}{杂交带}原本的宽度更窄：无法
形成空间\omterm{prop:b2:speciation:reinforcement}{渐变群}；杂交（若发生）影响
整个种群，只有选型交配才能使二者
保持分开。
\textbf{16.} 只与自身鸣声和喙形相同的雄性交配的雌性，
留下的孙辈比杂交的雌性更多；
在两个物种共同生活之处，偏好与线索（鸣声、喙形）
比单独生活之处更为分明。
\textbf{17.} $0.005\times 15\times 15 = 1.1$ 个不相容位点：平均
约为一个，杂种很可能能育性降低但并非不育。尚不
是生物学标准下的物种，正在形成之中。
\textbf{18.} 生物学标准：尚不是（杂种\omterm{def:b2:population-genetics:fitness}{适合度}几乎正常）。
形态学标准：或许是，若喙或羽色已出现
明显\omterm{def:b2:cell-differentiation:differentiation}{分化}。系统发生标准：是，若每个种群都具有固定的
诊断性差异。这几种概念恰恰在物种形成的
过程中产生分歧，这正是它们的价值所在。
\textbf{19.} $1/(5\times 10^{6}) = 2\times 10^{-7}$ 每物种每
年；对于 $10^{7}$ 个物种，每年两次灭绝。
\textbf{20.} 每年 200 至 2000 个物种。
\textbf{21.} $0.75\times 10^{7}/10^{5} = 75$ 个物种每年——低端
低于目前的速率。若当前事件
持续下去，便是一次速度超过五次\omterm{prop:b2:speciation:tempo}{大灭绝}的
大规模灭绝。
\textbf{22.} 失败者所空出的\omterm{prop:b2:speciation:displacement}{生态位}上
没有竞争者，因此任何能够利用其中一个的幸存变异
都会得到青睐，幸存者便\omterm{def:b2:cell-differentiation:differentiation}{分化}进入这些\omterm{prop:b2:speciation:displacement}{生态位}。
\textbf{23.} \omterm{prop:b2:speciation:tempo}{间断平衡}；异域（边缘）
物种形成发生于小而隔离的种群中，其后代
随后扩散进入主种群的分布范围。
\textbf{24.} $1.04^{n} = 2$：$n = \ln 2/\ln 1.04 = 18$ 代。
化石记录显示需要数百万年，是因为选择会发生逆转，
\omterm{prop:b2:population-genetics:kinds}{稳定选择}在大多数时间占主导，净变化
只是大量相反阶段之间的微小差值。
\textbf{25.} 置换之前：稳定共存，$N_1^{*} =
1151$，$N_2^{*} = 55$；之后：1176 和 47。喙的变化
\qty{-0.35}{mm}（一代内）。\omterm{prop:b2:speciation:reinforcement}{杂交带}宽度 \qty{3.2}{km}。
\end{solution}
